\subsection*{Discrete-Time Connectivity Model}

We formulate the epidemic model in discrete time with weekly time steps,
matching the natural resolution of surveillance data. Consider a system
of $n$ interconnected regions, where $S_i(t)$, $I_i(t)$, and $R_i(t)$
denote the fractions of Susceptible, Infected, and Recovered individuals
in region $i$ at time $t$. Since these are fractions, we have
$S_i(t) + I_i(t) + R_i(t) = 1$ for all $i$ and $t$.

Let $\mathbf{I}(t) = (I_1(t), \ldots, I_n(t))^T$ denote the vector of
infected fractions. The force of infection in region $i$ at time $t$ is:
\begin{equation*}
\lambda_i(t) = \beta_i(t) \sum_{j=1}^n C_{ij} I_j(t)
\end{equation*}
where $\mathbf{C} \in \mathbb{R}^{n \times n}$ is the contact matrix with
entries $C_{ij} \in [0,1]$ representing relative connectivity from
region $j$ to region $i$, and $\beta_i(t)$ is the region-specific
time-varying transmission rate with seasonal forcing:
\begin{equation*}
\beta_i(t) = \beta_0\left(1 + A\sin\left(\frac{2\pi t}{P} + \varphi_i\right)\right)
\end{equation*}
where $\beta_0$ is the baseline transmission rate, $A$ is the seasonal
amplitude, $P$ is the period, and $\varphi_i$ is the phase for region
$i$. Regions with identical phases ($\varphi_1 = \varphi_2$) experience
synchronized seasonal forcing, while different phases produce
asynchronous epidemic dynamics. We compare synchronized
($\varphi_1 = \varphi_2 = 0$) versus unsynchronized
($\varphi_1 = 0, \varphi_2 = \pi$) scenarios to assess how phase
alignment affects parameter identifiability.

Over one time step, the probability that a susceptible individual in
region $i$ does \emph{not} become infected is $\exp(-\lambda_i(t))$,
yielding the infection probability $p_i(t) = 1 - \exp(-\lambda_i(t))$.
The expected incidence (new infections) in region $i$ is:
\begin{equation}\label{eq:incidence}
\mu_i(t) = S_i(t) \left[1 - \exp(-\lambda_i(t))\right]
\end{equation}

Similarly, the probability an infected individual remains infected over
one time step is $\exp(-\gamma)$, where $\gamma$ is the recovery rate.
The discrete-time SIR dynamics are:
\begin{align}\begin{split}\label{eq:discrete_sir}
S_i(t+1) &= S_i(t) \exp(-\lambda_i(t)) \\
I_i(t+1) &= I_i(t) \exp(-\gamma) + \mu_i(t) \\
R_i(t+1) &= R_i(t) + I_i(t) [1 - \exp(-\gamma)]
\end{split}\end{align}

For our two-region model ($n=2$), we parameterize the contact matrix as:
\begin{equation}\label{eq:contact}
\mathbf{C} = \begin{bmatrix} 1-\theta & \theta \\ \theta & 1-\theta \end{bmatrix}
\end{equation}
where $\theta \in [0, 0.5]$ is the connectivity parameter representing
the fraction of contacts that occur between regions. When $\theta = 0$,
the regions are isolated; when $\theta = 0.5$, contacts are equally
distributed within and between regions. Row sums equal unity, reflecting
frequency-dependent transmission. We calibrate model parameters to
seasonal influenza following~\citet{keeling2008}; see
Table~\ref{tab:params}.

\begin{table}[h]
\centering
\begin{tabular}{lll}
\toprule
Parameter & Value & Description \\
\midrule
$R_0$ & 3.65 & Basic reproduction number \\
$\gamma$ & $7/2.2 \approx 3.18$ & Weekly recovery rate \\
$\beta_0$ & $\approx 2.06$ & Baseline transmission rate \\
$A$ & 0.7 & Seasonal amplitude \\
$P$ & 53 weeks & Seasonal period \\
$\rho$ & 0.3 & Reporting rate \\
%% $\theta$ & 0.05 & True connectivity (synthetic data) \\
\bottomrule
\end{tabular}
\caption{Model parameters calibrated to seasonal influenza.}
\label{tab:params}
\end{table}


\subsection*{Sensitivity Equations}

Computing the Cram\'er-Rao bound and optimizing the likelihood both
require derivatives of model outputs with respect to parameters. Let
$\phi \in \{\theta, S_1(0), S_2(0), I_1(0), I_2(0)\}$ denote an
arbitrary parameter. We derive the sensitivities by differentiating
the discrete dynamics.

The sensitivity of the force of infection is:
\begin{equation}\label{eq:sens_lambda}
\frac{\partial \lambda_i(t)}{\partial \phi} = \beta_i(t) \left(
\frac{\partial \mathbf{C}}{\partial \phi} \mathbf{I}(t) + \mathbf{C}
\frac{\partial \mathbf{I}(t)}{\partial \phi} \right)_i
\end{equation}
For the connectivity parameter:
\begin{equation*}
\frac{\partial \mathbf{C}}{\partial \theta} =
\begin{bmatrix} -1 & 1 \\ 1 & -1 \end{bmatrix}
\end{equation*}
while $\partial \mathbf{C}/\partial \phi = \mathbf{0}$ for $\phi \neq
\theta$.

The sensitivity of the expected incidence follows from
Eq.~\eqref{eq:incidence} by the product rule:
\begin{equation}\label{eq:sens_mu}
\frac{\partial \mu_i(t)}{\partial \phi} =
\frac{\partial S_i(t)}{\partial \phi} \left[1 - \exp(-\lambda_i(t))\right]
+ S_i(t) \exp(-\lambda_i(t)) \frac{\partial \lambda_i(t)}{\partial \phi}
\end{equation}

The state sensitivities propagate according to:
\begin{align}\begin{split}\label{eq:sens}
\frac{\partial S_i(t+1)}{\partial \phi} &= \exp(-\lambda_i(t))
\frac{\partial S_i(t)}{\partial \phi} - S_i(t) \exp(-\lambda_i(t))
\frac{\partial \lambda_i(t)}{\partial \phi} \\
\frac{\partial I_i(t+1)}{\partial \phi} &= \exp(-\gamma)
\frac{\partial I_i(t)}{\partial \phi} +
\frac{\partial \mu_i(t)}{\partial \phi}
\end{split}\end{align}

Initial conditions for the sensitivities are: $\partial
S_i(0)/\partial S_j(0) = \delta_{ij}$, $\partial I_i(0)/\partial
I_j(0) = \delta_{ij}$, and zero otherwise.  These equations are solved
jointly with the state equations in a single forward pass.


\subsection*{Observation Model}

We assume a constant reporting rate $\rho$, meaning a fraction $\rho$
of infections are observed. Each infection is reported independently
with probability $\rho$, so observations follow a binomial
distribution:
\begin{equation}\label{eq:obs_binomial}
Y_i(t) \sim \text{Binomial}(\mu_i(t), \rho)
\end{equation}
with mean $\mathbb{E}[Y_i(t)] = \rho \mu_i(t)$ and variance
$\text{Var}[Y_i(t)] = \rho(1-\rho) \mu_i(t)$.

For analytical tractability, we approximate this by a Gaussian
distribution via the central limit theorem:
\begin{equation}\label{eq:obs_gaussian}
Y_i(t) = \rho \mu_i(t) + \epsilon_i(t), \quad
\epsilon_i(t) \sim \mathcal{N}(0, \sigma_i^2(t))
\end{equation}
where $\sigma_i^2(t) = \rho(1-\rho)\mu_i(t)$ is the observation variance.
This approximation is accurate when $\mu_i(t)$ is large.

Under this Gaussian model, the likelihood of observing data
$\mathbf{Y} = \{Y_i(t)\}$ given parameters $\boldsymbol{\phi}$ is:
\begin{equation}\label{eq:likelihood}
\mathcal{L}(\boldsymbol{\phi}) = \prod_{t=0}^{T-1} \prod_{i=1}^{n}
\frac{1}{\sqrt{2\pi\sigma_i^2(t)}}
\exp\left(-\frac{(Y_i(t) - \rho\mu_i(t))^2}{2\sigma_i^2(t)}\right)
\end{equation}

The log-likelihood is:
\begin{equation}\label{eq:loglik}
\ell(\boldsymbol{\phi}) = -\frac{1}{2}\sum_{t=0}^{T-1}\sum_{i=1}^{n}
\left[\log(2\pi \sigma_i^2(t)) +
\frac{(Y_i(t) - \rho\mu_i(t))^2}{\sigma_i^2(t)}\right]
\end{equation}



\subsection*{CRLB for the Discrete SIR Model}

We now derive the Fisher Information Matrix for our Gaussian observation
model (Eq.~\ref{eq:obs_gaussian}). Taking the partial derivative of the
log-likelihood (Eq.~\ref{eq:loglik}) with respect to parameter $\phi$:
\begin{equation*}
\frac{\partial \ell}{\partial \phi} = -\frac{1}{2}\sum_{t,i} \left[
\frac{1}{\mu_i(t)} - \frac{2\rho(Y_i(t) - \rho\mu_i(t))}{\rho(1-\rho)\mu_i(t)}
- \frac{(Y_i(t) - \rho\mu_i(t))^2}{\rho(1-\rho)\mu_i(t)^2}\right]
\frac{\partial \mu_i(t)}{\partial \phi}
\end{equation*}

This can be written as:
\begin{equation}\label{eq:score}
\frac{\partial \ell}{\partial \phi} = \sum_{t,i} A_i(t)
\frac{\partial \mu_i(t)}{\partial \phi}
\end{equation}
where
\begin{equation*}
A_i(t) = -\frac{1}{2\mu_i(t)} + \frac{Y_i(t) - \rho\mu_i(t)}{(1-\rho)\mu_i(t)}
+ \frac{(Y_i(t) - \rho\mu_i(t))^2}{2\rho(1-\rho)\mu_i(t)^2}
\end{equation*}

The Fisher Information Matrix is:
\begin{equation*}
J_{kl} = \mathbb{E}\left[\frac{\partial \ell}{\partial \phi}
\frac{\partial \ell}{\partial \phi_l}\right] = \sum_{t,i}\sum_{s,j}
\mathbb{E}[A_i(t) A_j(s)] \frac{\partial \mu_i(t)}{\partial \phi_k}
\frac{\partial \mu_j(s)}{\partial \phi_l}
\end{equation*}

We can easily verify that $\mathbb{E}[A_i(t)]]= 0$. Since observations
  are independent across time and regions, $\mathbb{E}[A_i(t) A_j(s)]=
  0$ when $(i,t) \neq (j,s)$. For the diagonal terms, we need
  $\mathbb{E}[A_i(t)^2]$. Writing $\epsilon_i(t) = Y_i(t) -
  \rho\mu_i(t)$:
\begin{equation*}
A_i(t) = -\frac{1}{2\mu_i(t)} + \frac{\epsilon_i(t)}{(1-\rho)\mu_i(t)}
+ \frac{\epsilon_i(t)^2}{2\rho(1-\rho)\mu_i(t)^2}
\end{equation*}

Expanding $A_i(t)^2$:
\begin{align*}
A_i(t)^2 &= \frac{1}{4\mu_i(t)^2} - \frac{\epsilon_i(t)}{(1-\rho)\mu_i(t)^2}
- \frac{\epsilon_i(t)^2}{2\rho(1-\rho)\mu_i(t)^3} \\
&\quad + \frac{\epsilon_i(t)^2}{(1-\rho)^2\mu_i(t)^2}
+ \frac{\epsilon_i(t)^3}{\rho(1-\rho)^2\mu_i(t)^3}
+ \frac{\epsilon_i(t)^4}{4\rho^2(1-\rho)^2\mu_i(t)^4}
\end{align*}

Using the moments of $\epsilon_i(t) \sim \mathcal{N}(0, \rho(1-\rho)\mu_i(t))$:

\begin{align*}
  \mathbb{E}[\epsilon] &= 0, \\
  \mathbb{E}[\epsilon^2] &= \rho(1-\rho)\mu_i(t), \\
  \mathbb{E}[\epsilon^3] &= 0, \text{ and } \\
  \mathbb{E}[\epsilon^4] &= 3[\rho(1-\rho)\mu_i(t)]^2.
\end{align*}
  Taking expectations:
\begin{align*}
\mathbb{E}[A_i(t)^2] &= \frac{1}{4\mu_i(t)^2}
- \frac{\rho(1-\rho)\mu_i(t)}{2\rho(1-\rho)\mu_i(t)^3}
+ \frac{\rho(1-\rho)\mu_i(t)}{(1-\rho)^2\mu_i(t)^2}
+ \frac{3[\rho(1-\rho)\mu_i(t)]^2}{4\rho^2(1-\rho)^2\mu_i(t)^4} \\
&= \frac{1}{4\mu_i(t)^2} - \frac{1}{2\mu_i(t)^2}
+ \frac{\rho}{(1-\rho)\mu_i(t)^2} + \frac{3}{4\mu_i(t)^2} \\
&= \frac{1}{2\mu_i(t)^2} + \frac{\rho}{(1-\rho)\mu_i(t)^2}
= \frac{1+\rho}{2(1-\rho)\mu_i(t)^2}
\end{align*}

Therefore:
\begin{equation}\label{eq:fim_scalar}
J_{kl} = \frac{1+\rho}{2(1-\rho)}\sum_{t=0}^{T-1}\sum_{i=1}^{n}
\frac{1}{\mu_i(t)^2} \frac{\partial \mu_i(t)}{\partial \phi}
\frac{\partial \mu_i(t)}{\partial \phi_l}
\end{equation}

In matrix form, define the Jacobian $\mathbf{G} \in \mathbb{R}^{2T \times p}$
with entries $G_{(t,i),k} = \partial \mu_i(t)/\partial \phi_k$---the
incidence sensitivities from Eq.~\eqref{eq:sens_mu}---and the diagonal
weight matrix:
\begin{equation*}
\mathbf{W} = \frac{1+\rho}{2(1-\rho)} \,
\text{diag}\left(\mu_1(0)^{-2}, \ldots, \mu_n(T-1)^{-2}\right)
\end{equation*}
Then:
\begin{equation}\label{eq:fim}
\mathbf{J} = \mathbf{G}^T \mathbf{W} \mathbf{G}
\end{equation}
and we find the CRLB as $\text{Var}(\hat{\theta}) \geq
[\mathbf{J}^{-1}]_{11}$.
%% The CRLB for connectivity is $[\mathbf{J}^{-1}]_{11}$, computed via
%% Cholesky factorization: if $\mathbf{J} = \mathbf{L}\mathbf{L}^T$, then
%% $[\mathbf{J}^{-1}]_{11} = \|\mathbf{L}^{-1}\mathbf{e}_1\|^2$ where
%% $\mathbf{e}_1$ is the first standard basis vector.

\subsubsection*{Statistical Power}

The CRLB determines the statistical power for detecting non-zero
connectivity. Specifically, we test
\begin{align*}
  H_0&: \theta = 0\\
  H_1&: \theta > 0
\end{align*}

at significance level $\alpha$. We assume our estimator $\hat{\theta}$
is both efficient and unbiased, hence it achieves the CRLB and
$\hat{\theta} \sim \mathcal{N}(\theta, CRLB)$. We thus reject the null
$H_0$ whenever $\hat{\theta} > z_\alpha\sqrt{\text{CRLB}}$. The
statistical power is the probability of correctly rejecting $H_0$ i.e.
\begin{align*}%\label{eq:power}
  \text{Power} &= \Pr(\hat{\theta} > \sqrt{\text{CRLB}} z_\alpha|H_1) \\
  &= \Pr(\frac{\hat{\theta} - \theta}{\sqrt{\text{CRLB}}} > z_\alpha|H_1) \\
  &= \Pr(Z > z_\alpha - \frac{\theta}{\sqrt{\text{CRLB}}}) \\
  &= 1 - \Phi\left(z_\alpha -
  \frac{\theta}{\sqrt{\text{CRLB}}}\right)
\end{align*}
where $\Phi$ is the standard normal CDF and $z_\alpha$ is the critical
value ($\Pr(Z < z_{0.05}) = \alpha$ (e.g.~ $z_{0.05}= 1.645$).

%% Power approaches 1 when the $\theta$ is large relative to the
%% estimation uncertainty $\sqrt{\text{CRLB}}$, and approaches
%% $\alpha$ when this ratio is small.

